% ---- pg-443 （书页 431） ----

\noindent\textbf{显示板 10.1（续）}

\begin{mvbox}
{\ttfamily\small
Canonical Correlation Analysis\par
\vspace{8pt}
{\centering
\begin{tabular}{lllll}
 & Canonical & Adjusted & Approx & Squared\\
 & Correlation & Canonical & Standard & Canonical\\
 & & Correlation & Error & Correlation\\[4pt]
1 & 0.631 085 & 0.628 291 & 0.036 286 & 0.398 268\\
2 & 0.056 794 & & 0.060 108 & 0.003 226
\end{tabular}\par}
\vspace{12pt}

{\centering
\fbox{Raw Canonical Coefficient for the 'VAR' Variables}\hfill{\bfseries\rmfamily OUTPUT}\par}
\vspace{8pt}
{\centering
\begin{tabular}{lll}
 & HEAD1 & HEAD2\\
X1 & 0.780 792 438 9 & $-0.855$ 973 184\\
X2 & 0.344 506 830 1 & 1.106 183 514 5
\end{tabular}\par}
\vspace{12pt}

{\centering
\fbox{Raw Canonical Coefficient for the 'WITH' Variables}\par}
\vspace{8pt}
{\centering
\begin{tabular}{lll}
 & LEG1 & LEG2\\
X3 & 0.060 250 877 5 & $-2.648$ 156 338\\
X4 & 0.943 948 961 & 2.474 938 891 3
\end{tabular}\par}
\vspace{12pt}

Canonical Structure\par
\vspace{8pt}
{\centering
\fbox{Correlations Between the 'VAR' Variables and Their Canonical Variables}\par}
\vspace{8pt}
{\centering
\begin{tabular}{llll}
 & HEAD1 & HEAD2 & (see 10-29)\\
X1 & 0.954 8 & $-0.297$ 4 &\\
X2 & 0.738 8 & 0.673 9 &
\end{tabular}\par}
\vspace{12pt}

{\centering
\fbox{Correlations Between the 'WITH' Variables and Their Canonical Variables}\par}
\vspace{8pt}
{\centering
\begin{tabular}{llll}
 & LEG1 & LEG2 & (see 10-29)\\
X3 & 0.934 3 & $-0.356$ 4 &\\
X4 & 0.999 7 & 0.022 7 &
\end{tabular}\par}
\vspace{12pt}

{\centering
\fbox{\begin{tabular}{c}Correlations Between the 'VAR' Variables\\ and the Canonical Variables of the 'WITH' Variables\end{tabular}}\par}
\vspace{8pt}
{\centering
\begin{tabular}{llll}
 & LEG1 & LEG2 & (see 10-29)\\
X1 & 0.602 5 & $-0.016$ 9 &\\
X2 & 0.466 3 & 0.038 3 &
\end{tabular}\par}
\vspace{12pt}

{\centering
\fbox{\begin{tabular}{c}Correlations Between the 'WITH' Variables\\ and the Canonical Variables of the 'VAR' Variables\end{tabular}}\par}
\vspace{8pt}
{\centering
\begin{tabular}{llll}
 & HEAD1 & HEAD2 & (see 10-29)\\
X3 & 0.589 7 & $-0.020$ 2 &\\
X4 & 0.630 9 & 0.001 3 &
\end{tabular}\par}
\par}
\end{mvbox}

\noindent\textbf{例 10.5}\quad (职业满意度的典型相关分析)

在研究组织结构对“职业满意度”的影响时, 作为其中一部分, 邓纳姆(Dunham, 见文献[4])调查了职业满意度与职业特性相关的程度. 利用一调查方法, 对从一大型零售公司各分公司
